% Chapter 1
\chapter{مقدمه}

%\section{مقدمه}

در حال حاضر، استفاده از دوقلوی دیجیتال\LTRfootnote{Digital Twins} در صنایع مختلف، از جمله تولید، بهداشت و درمان، حمل و نقل و حتی شهرسازی، در حال افزایش است. این تکنولوژی به سازمان‌ها این امکان را می‌دهد که فرآیندها و محصولات خود را به شکل مؤثری مدیریت کنند و تصمیمات هوشمندانه‌ای اتخاذ نمایند. به طور خاص، توانایی پیش‌بینی خرابی‌ها و مشکلات پیش‌رو از طریق تجزیه و تحلیل داده‌های جمع‌آوری شده، به شرکت‌ها، کارخانه‌ها و مدیران کمک می‌کند تا به سمت کاهش هزینه‌های عملیاتی و افزایش بهره‌وری حرکت کنند.


تاریخچه استفاده از شبیه‌ساز‌های دیجیتال به سال‌ها قبل برمی‌گردد. مشهورترین مثال از این کاربردها، سیستم شبیه‌سازی که در طی ماموریت آپولوی 13 توسط ناسا ایجاد شد، است. در این ماموریت، زمانی که یکی از مخازن اکسیژن در فضاپیما منفجر شد، تیم کنترل پرواز با استفاده از یک مدل شبیه‌سازی شده از فضاپیما و اجزای آن، توانستند راه‌حلی برای عبور از این مشکل پیدا کنند. این واقعه نشان داد که چطور شبیه‌سازی‌ها می‌توانند به مهندسان کمک کنند تا از راه دور روندهای پیچیده را در محیطی امن آزمایش کنند و چالش‌های زندگی واقعی را شبیه‌سازی کنند. هرچند این شبیه‌سازی در آن زمان، به‌عنوان دوقلوهای دیجیتال شناخته نمی‌شد، اما تأثیر آنها در درک عمیق‌تر از ارتباط میان فضای فیزیکی و دیجیتال غیرقابل انکار است.
\cite{b3}

در نهایت، با پیشرفت مستمر فناوری‌های دیجیتال، عصر جدیدی از نوآوری و کارایی در انتظار صنایع مختلف است. دوقلوی دیجیتال به عنوان یکی از مؤلفه‌های کلیدی این تحول، نقش مهمی در ارتباط و همگرایی میان دنیای فیزیکی و دیجیتال ایفا خواهد کرد. بدین ترتیب، می‌توان انتظار داشت که کاربردهای دوقلوی دیجیتال در آینده نزدیک گسترش یافته و به بهبود مستمر کیفیت خدمات و محصولات منجر شود.

در این پژوهش پس از اینکه به تعریف چنین مفاهیم مختلف در این زمینه پرداخته می‌شود، به کاربرد آنها در زمینه‌های مختلف اشاره شده و مثال‌هایی از هرکدام بیان می‌شود. سپس به بیان یک معماری پیشنهادی برای دوقلوی دیجیتال چندعاملی پرداخته شده و جزییات آن شرح داده می‌شود. در نهایت چند ابزار مناسب در استفاده از دوقلوی دیجیتال معرفی می‌گردد.


